\section{Consequences and comparison with existing theory}\label{sec:scope}
The central theorem separates actual acquisition, geometry and causal implementation, then proves a quantitative connection. Its forward gain matrix need not be diagonal in the predictive strata. Its occupation comparison is derived from raw invariant-reference bounds in Proposition~\ref{prop:nr-reference}, without identifying the reference measure with the actual law. Lemma~\ref{lem:nr-overlap} permits a finite multiplicity of overlapping score neighborhoods. Proposition~\ref{prop:nr-acquisition} treats atomic acquired laws by retaining the unresolved-variance term rather than imposing a false fine-scale small-ball bound.

The recurrent prefix experiment gives an explicit common-task law $\delta^2+\E\ell_{\min(H_T,\lfloor\log_2M\rfloor)}^2$. Its known depth is generated by the very same nonreset acquisition process that the encoder must compress. The older one-probe and selected-menu minimax theorems remain independent sharpness witnesses, not proofs of adaptive acquisition or of a universal full-resource region. The exact and noisy continuation converses likewise constrain information that truly survives a cut, rather than every register in one preferred implementation.

\subsection{Theorem-level comparison}
\paragraph{Predictive states and sufficient statistics.}
Shalizi and Crutchfield's minimality and uniqueness theorems for causal states \cite[Theorems 2--3]{shalizi} identify minimal prescient representations of a stochastic process. Predictive state representations \cite{psr} use future action--observation tests. Our quotient proposition does not claim priority for either construction. It explicitly includes controlled continuation legality, paid resources and conditional-version domains. Its extra measurable work is the density-instrument descent at every declared continuous report. The categorical treatment of kernels and sufficiency in \cite{fritz} supplies another general language; an abstract sufficient statistic alone does not provide the acquired-mass or finite-machine estimates used here.

\paragraph{Bisimulation and belief-state control.}
Behavioral metrics for Markov decision processes \cite{ferns} compare states through rewards and transition laws and control value differences. Our metric is fixed by executable prediction tasks and valid update geometry; it is not asserted to be a universal bisimulation metric. Saldi, Y\"uksel and Linder \cite[Assumption 1 and Theorem 3]{saldi} obtain near-optimal finite belief-model policies for discounted partially observed control under their continuity, compactness and moment hypotheses. That theorem optimizes a controller, whereas our matching lower is for a fixed exploration and the law it actually acquires. Our finite-horizon two-sided resolution estimate does not imply their unrestricted control conclusion, nor does asymptotic policy near-optimality supply our actual-mass lower bound.

\paragraph{Quantization and nonlinear filtering.}
Conditional-mean projection and optimal probability quantization are classical \cite{graf}. Nonlinear filter approximation by quantization is developed in \cite{pages,sellami,feuer}. The additional obligation here is a matching lower for the very same acquired task, together with the forced-moment condition needed by a persistent finite machine. Zhu's local estimates for Ahlfors--David measures \cite{zhu}, and for Moran measures \cite{zhu-moran}, concern fine uniformity of contributions of optimal cells and successive quantization errors. The cylinder covers in Lemma~\ref{lem:moran-profile} and Theorem~\ref{thm:nr-moran} are not offered as replacements for that theory. Its role is to verify a multiscale profile, then propagate that profile through an actual nonstationary recurrent experiment with countably many weighted components. The finite-budget anisotropic results additionally retain width collapse and finite intersecting mixtures. We do not classify arbitrary singular measures.

\paragraph{Positive moment operators.}
Average contraction of iterated random functions \cite{diaconis} and moment stability of Markov jump systems \cite{costa} are established subjects. Ogura and Martin \cite{ogura} characterize mean stability for positive semi-Markovian jump systems with resets by a spectral-radius condition. That homogeneous stability question is not the joint domination of forced error measures by changing acquired measures required here. We claim no novelty for a Neumann resolvent or the criterion $\rho(K_2)<1$. The role of Theorem~\ref{thm:kernel} is more specific: reverse-edge operators propagate the conditional first and mixed second moments of the \emph{rounding forcing}, and the resulting matrix is normalized by the actual acquisition weights. Exact suspension cancels in the forced profiles, not in the physical budget. This distinction is why a recurrent expanding kernel can still yield a matching finite-memory risk curve. The new forward drift in Theorem~\ref{thm:nonreset} controls all cross-stratum old-error edges and all allocation-dependent radii, after occupation comparison, without that covariance normalization. Its matrix condition is standard positive stability, but its coupling to the actual-law lower and executable online construction is the claim studied here. In the complementary reset-preserved class, Theorem~\ref{thm:bridge} propagates forward error measures and obtains the terminal mass factor directly; Proposition~\ref{prop:recharge-flow} identifies the necessary refresh-versus-expansion balance used by that sufficient criterion. We do not claim that this sufficient condition is a necessary spectral classification of all raw kernels.

\paragraph{Causal rate distortion and sequential reconstruction.}
Nonanticipative rate-distortion theory imposes realizability constraints on reconstruction kernels and relates information minimization to filtering on abstract spaces \cite{nonanticipative}. A communication or directed-information budget is not, without a separate coding theorem, the cardinality of the complete state retained across each computational cut. Our theorem gives explicit recursive labels and a checkpoint converse in the same task, not an identification of those resource models. The finite-prefix acquisition theorem counts word length and mode in a fused finite alphabet, and separately records input, arithmetic and controller costs.

\paragraph{Comparison of experiments.}
Blackwell comparison \cite{blackwell}, Le Cam's deficiency theory \cite{lecam}, and Torgersen's systematic account \cite{torgersen} provide the classical common-decision-loss principle. The triangle inequality and total-variation contraction in Theorem~\ref{thm:morphism} are instances of that principle. The added structure is the finite-resource causal implementation, with explicit lifted controller classes, stopping indices, clocks and numerical interfaces. A forward simulator yields an upper risk transfer, not a lower bound relative to a different oracle. The mandatory-state and robust converses are proved directly in Theorems~\ref{thm:tag} and \ref{thm:robust}, rather than attributed to a forward deficiency bound.

\subsection{Necessary distinctions illustrated by exact examples}
First, the valid envelope $[0,1]$ equipped with a point mass and the same envelope equipped with a uniform law have respectively zero one-point distortion and strictly positive finite-label distortion. Hence a reachability envelope is not an acquired law. More generally a component of probability $w$ contributes with weight $w$, even if its geometric diameter is large.

Second, an identity update has no need for repeated quantization. If the initial state has already been replaced by a representative, the same representative can be retained indefinitely. The bound $Nr$ obtained by blindly injecting a covering radius at every hold is not an inherent resource requirement. The acquisition-weighted certificate corrects precisely this failure of accounting.

Third, a statistical garbling can increase memory needed to match its \emph{own} oracle. Observing a fair hidden bit exactly produces only posterior probabilities zero and one. Adding independent Gaussian noise to this observation produces a continuously varying posterior probability. Two labels reproduce the first oracle exactly, but no finite number reproduces the latter continuously distributed oracle exactly under square score. This is compatible with data processing because the baselines differ.

Fourth, let $P_N=\delta_0$ and $Q_N=(1-e^{-Nc})\delta_0+e^{-Nc}\delta_1$, $c>0$. Their total variation distance tends to zero, but the exponential rates at the point one are infinite and $c$, respectively. The common-risk defect of this paper is therefore not a certificate of large-deviation equivalence.

Finally, a hidden mode or an internal simulator register is not free because its name is omitted from a state vector. Independent pairs of retained states have the product cardinality in Theorem~\ref{thm:morphism}. Conversely that product is a constructive upper bound, not a theorem that every simulator requires it. Theorem~\ref{thm:tag} establishes exact support necessity through zero-error continuation challenges; Theorem~\ref{thm:soft-state} replaces it by an information lower bound when errors are allowed. An exact label machine and a finite-bit numerical implementation are also different objects: the latter retains the nonzero calibration, workspace, and readout terms of \eqref{eq:joint}. Exact-arithmetic zero distortion is not a promise of a finite-precision physical output with zero error.


\subsection{Uniformity and the remaining mathematical questions}
Theorem~\ref{thm:nonreset} allows nonreset transfer between strata, expanding edges, nonstationary acquired laws comparable to an invariant reference, and finite or countable geometric partitions with bounded overlap. The uniform constants are explicit in the drift and occupation bounds. They are not an assertion that every raw kernel admits such bounds. In particular, rapid creation or extinction of strata can violate occupation comparison; the separate reset-preserved forward-measure theorem still handles its own nonstationary refresh-flow class. No necessity characterization of optimal finite-memory prediction is inferred from the spectral-radius criterion for a comparison matrix.

The countable overlap result requires finite multiplicity at the declared relative scales and dyadic regularity of the profiles. Arbitrary unbounded-multiplicity intersections or all singular measures remain outside it. The finite intersecting-stratum and rank-collapse results remain independently available. The singular nonreset examples prove a nonuniform recurrent extension, not a classification of all nonlinear filters, diffusions or hidden-state systems.

The recurrent acquired-depth law treats a known raw kernel and a fixed, fully charged acquisition protocol, with a two-point calibration ambiguity that the reports cannot learn. It does not optimize the division of the raw budget between initial readings and recurrence, estimate unknown transition kernels, or supply an optimal pilot--calibration design. These remain mathematical responsibilities within the same mother problem. A universal matching region for program length, transient workspace, physical runtime and simulator state is also not proved. The information lower bounds and explicit finite-bit implementations give narrower, stated claims rather than an equality for that entire vector.

Finally, a deterministic-horizon energy estimate is not automatically a stopping-time estimate: a stopping rule can select large errors. The bounded stopped common-law comparisons of Theorem~\ref{thm:morphism} retain their separate uniform hypotheses. Infinite-horizon transcript comparison, large-deviation equivalence and unbounded-generator closure require different estimates. Neither successful compilation nor finite regression checks establish any of the continuum theorems proved here.
